\section{Arithmetic obstructions from the complement quotient}
\label{sec:arithmetic}

The uniform low-root bound rules out common exponent divisors greater
than two. Reduction of the characteristic polynomial gives further
infinite congruence classes, even when another quotient root is an integer.

\begin{theorem}
\label{thm:common-divisor}
Every positive three-prime exponent triple with $\gcd(a,b,c)\geq3$
gives a nonintegral ideal intersection graph.
\end{theorem}

\begin{proof}
Put $d=\gcd(a,b,c)$. All entries of the six-support complement quotient
$C$ are integer combinations of the exponents and their pairwise
products, so $C/d$ is an integer matrix. Every rational eigenvalue of
an integer matrix is an integer, by the rational-root theorem for its
monic characteristic polynomial. Therefore every integer eigenvalue
of $C$ is a multiple of $d$.

Order $a\leq b\leq c$. If $a=3$, Theorem~\ref{thm:minimum-three}
applies. Otherwise $a\geq4$, and Theorem~\ref{thm:uniform-low-root}
supplies an eigenvalue $0<\mu<3$. No positive multiple of $d\geq3$
lies in this interval. Hence $\mu$ and its graph lift $V-\mu$ are
noninteger.
\end{proof}

In particular, if $m$ has exactly three distinct prime factors, then
$G_{m^d}$ is not Laplacian integral for every $d\geq3$.

\begin{theorem}
\label{thm:congruence}
A positive three-prime exponent triple gives a nonintegral ideal
intersection graph under either of the following conditions:
\begin{enumerate}
\item all three exponents are odd;
\item there is an odd prime $\ell$ and a nonzero residue $t$ such that
$a\equiv b\equiv c\equiv t\pmod\ell$ and $t^2+8t+4$ is a quadratic
nonresidue modulo $\ell$.
\end{enumerate}
In particular, common residues $1,3,4$ modulo $5$, or $1,2,4,5$ modulo
$7$, imply nonintegrality.
\end{theorem}

\begin{proof}
The monic integer quintic $h_C(x)=\det(xI-C)/x$ would split into
linear factors modulo every prime if all its roots were integers.
A nonlinear irreducible factor in one reduction therefore supplies a
noninteger real complement root and a noninteger graph lift.

If the exponents are odd, reduce the general quotient modulo two:
\[
h_C(x)\equiv x(x^2+x+1)^2\pmod2.
\]
The quadratic has no root in $\mathbb F_2$, proving the first assertion.
For a common residue $t$, the equal-exponent identity gives
\[
h_C(x)\equiv(x-t^2-t)
\bigl(x^2-(t^2+4t)x+3t^2\bigr)^2\pmod\ell.
\]
The quadratic discriminant is $t^2(t^2+8t+4)$. For odd $\ell$ and
nonzero $t$, the assumed nonresidue makes the quadratic irreducible.
Modulo five the discriminant factors for $t=1,3,4$ are $3,2,2$,
respectively; the squares are $0,1,4$. Modulo seven the corresponding
values for $t=1,2,4,5$ are $6,3,3,6$, outside the square residues.
\end{proof}

These conditions allow unequal exponents with arbitrarily large ratios
or differences. For example, $(11,16,21)$, $(13,18,23)$ and $(14,19,24)$
have gcd one and mixed parity, but meet the modulo-five criterion.
The other parity reductions are $x^5$ for three even exponents and
$x^2(x+1)^3$ for one or two odd exponents. Their complete splitting
does not establish integrality.

The arithmetic arguments do not assume that the smallest positive
quotient root is noninteger. The necessary endpoint-zero condition of
Corollary~\ref{cor:balanced-interval} remains insufficient: the already
settled boundary triple $(9,9,136)$ has $h_C(1)=0$ and other noninteger
quotient roots. The remaining three-prime problem has gcd at most two,
at least one even exponent, and lies outside the congruence classes
above as well as the earlier proved regions.
